% preamble:
% \usepackage{tikz}
% \usetikzlibrary{arrows.meta}
\section{2026-05-16}

\subsection{Boundary Points and Manifold Boundary}

\begin{center}
\begin{tikzpicture}[
  >=Stealth,
  every node/.style={font=\small},
  map/.style={->, thick},
  tmap/.style={->, thick, orange!80!black},
  outline/.style={thick},
  scale=1
]

% =====================================================
% Top: manifold with boundary M
% =====================================================

\node[anchor=west] at (-7.2,7.3)
  {\underline{Recap}: manifolds with boundary, transition map};

% M as a half-ellipse
\draw[outline]
  (-4.2,4.2) -- (4.2,4.2)
  arc[start angle=0,end angle=180,x radius=4.2,y radius=1.7]
  -- cycle;

\node at (0,6.05) {$M$};
\node at (3.15,4.00) {$\partial M$};

% exact shading of U \cap V
\begin{scope}
  % U half-ellipse
  \clip
    (-2.15,4.2) -- (1.05,4.2)
    arc[start angle=0,end angle=180,x radius=1.6,y radius=1.25]
    -- cycle;
  % V half-ellipse
  \clip
    (-1.05,4.2) -- (2.15,4.2)
    arc[start angle=0,end angle=180,x radius=1.6,y radius=1.25]
    -- cycle;
  \fill[orange!45, opacity=0.85] (-3,4.0) rectangle (3,5.8);
\end{scope}

% exact shading of U \cap V \cap \partial M
% Since this is lower-dimensional, we draw it as a thick red segment on the boundary.
\begin{scope}
  \clip
    (-2.15,4.2) -- (1.05,4.2)
    arc[start angle=0,end angle=180,x radius=1.6,y radius=1.25]
    -- cycle;
  \clip
    (-1.05,4.2) -- (2.15,4.2)
    arc[start angle=0,end angle=180,x radius=1.6,y radius=1.25]
    -- cycle;
  \draw[red!65, line width=5pt, opacity=0.75] (-4.2,4.2) -- (4.2,4.2);
\end{scope}

% Draw U and V after shading
\draw[blue!70!black, thick]
  (-2.15,4.2) -- (1.05,4.2)
  arc[start angle=0,end angle=180,x radius=1.6,y radius=1.25]
  -- cycle;

\draw[red!70!black, thick]
  (-1.05,4.2) -- (2.15,4.2)
  arc[start angle=0,end angle=180,x radius=1.6,y radius=1.25]
  -- cycle;

\node[blue!70!black] at (-1.55,5.25) {$U$};
\node[red!70!black] at (1.55,5.25) {$V$};

\node[orange!80!black] at (0,5.15) {$U\cap V$};
\node[red!70!black] at (1.35,3.78) {$U\cap V\cap\partial M$};

% point p on the boundary intersection
\fill[green!60!black] (0,4.2) circle (1.8pt);
\node[green!50!black, below] at (0,4.12) {$p$};

% =====================================================
% Left: phi-chart in H^n
% =====================================================

% H^n box
\draw[blue!70!black, thick] (-7.0,0.0) rectangle (-2.6,2.55);
\draw[blue!70!black, thick] (-7.0,0.0) -- (-2.6,0.0);
\node[blue!70!black] at (-6.65,2.32) {$\mathbb H^n$};
\node[blue!70!black] at (-3.10,-0.28) {$\partial\mathbb H^n$};

% phi(U): half-ellipse meeting boundary
\draw[blue!70!black, thick]
  (-6.15,0.0) -- (-3.45,0.0)
  arc[start angle=0,end angle=180,x radius=1.35,y radius=0.85]
  -- cycle;

\node[blue!70!black] at (-4.80,1.25) {$\varphi(U)$};

% phi(U cap V): smaller half-ellipse meeting boundary
\draw[orange!80!black, thick]
  (-5.05,0.0) -- (-3.75,0.0)
  arc[start angle=0,end angle=180,x radius=0.65,y radius=0.50]
  -- cycle;

% strict shading of phi(U cap V)
\begin{scope}
  \clip
    (-6.15,0.0) -- (-3.45,0.0)
    arc[start angle=0,end angle=180,x radius=1.35,y radius=0.85]
    -- cycle;
  \clip
    (-5.05,0.0) -- (-3.75,0.0)
    arc[start angle=0,end angle=180,x radius=0.65,y radius=0.50]
    -- cycle;
  \fill[orange!45, opacity=0.85] (-6.2,-0.1) rectangle (-3.4,1.1);
\end{scope}

% boundary part in phi(U cap V)
\begin{scope}
  \clip
    (-5.05,0.0) -- (-3.75,0.0)
    arc[start angle=0,end angle=180,x radius=0.65,y radius=0.50]
    -- cycle;
  \draw[red!65, line width=4pt, opacity=0.75] (-6,0) -- (-3,0);
\end{scope}

% redraw outlines
\draw[blue!70!black, thick]
  (-6.15,0.0) -- (-3.45,0.0)
  arc[start angle=0,end angle=180,x radius=1.35,y radius=0.85]
  -- cycle;

\draw[orange!80!black, thick]
  (-5.05,0.0) -- (-3.75,0.0)
  arc[start angle=0,end angle=180,x radius=0.65,y radius=0.50]
  -- cycle;

\fill[green!60!black] (-4.40,0.0) circle (1.6pt);
\node[green!50!black, below] at (-4.40,-0.05) {$\varphi(p)$};

\node[orange!80!black] at (-4.40,-0.58) {$\varphi(U\cap V)$};

% =====================================================
% Right: psi-chart in H^n
% =====================================================

% H^n box
\draw[red!70!black, thick] (2.6,0.0) rectangle (7.0,2.55);
\draw[red!70!black, thick] (2.6,0.0) -- (7.0,0.0);
\node[red!70!black] at (6.65,2.32) {$\mathbb H^n$};
\node[red!70!black] at (6.20,-0.28) {$\partial\mathbb H^n$};

% psi(V): half-ellipse meeting boundary
\draw[red!70!black, thick]
  (3.45,0.0) -- (6.15,0.0)
  arc[start angle=0,end angle=180,x radius=1.35,y radius=0.85]
  -- cycle;

\node[red!70!black] at (4.80,1.25) {$\psi(V)$};

% psi(U cap V): smaller half-ellipse meeting boundary
\draw[orange!80!black, thick]
  (3.75,0.0) -- (5.05,0.0)
  arc[start angle=0,end angle=180,x radius=0.65,y radius=0.50]
  -- cycle;

% strict shading of psi(U cap V)
\begin{scope}
  \clip
    (3.45,0.0) -- (6.15,0.0)
    arc[start angle=0,end angle=180,x radius=1.35,y radius=0.85]
    -- cycle;
  \clip
    (3.75,0.0) -- (5.05,0.0)
    arc[start angle=0,end angle=180,x radius=0.65,y radius=0.50]
    -- cycle;
  \fill[orange!45, opacity=0.85] (3.4,-0.1) rectangle (6.2,1.1);
\end{scope}

% boundary part in psi(U cap V)
\begin{scope}
  \clip
    (3.75,0.0) -- (5.05,0.0)
    arc[start angle=0,end angle=180,x radius=0.65,y radius=0.50]
    -- cycle;
  \draw[red!65, line width=4pt, opacity=0.75] (3,0) -- (6,0);
\end{scope}

% redraw outlines
\draw[red!70!black, thick]
  (3.45,0.0) -- (6.15,0.0)
  arc[start angle=0,end angle=180,x radius=1.35,y radius=0.85]
  -- cycle;

\draw[orange!80!black, thick]
  (3.75,0.0) -- (5.05,0.0)
  arc[start angle=0,end angle=180,x radius=0.65,y radius=0.50]
  -- cycle;

\fill[green!60!black] (4.40,0.0) circle (1.6pt);
\node[green!50!black, below] at (4.40,-0.05) {$\psi(p)$};

\node[orange!80!black] at (4.40,-0.58) {$\psi(U\cap V)$};

% =====================================================
% Chart maps and transition map
% =====================================================

\draw[map, blue!70!black]
  (-1.25,4.35) to[out=220,in=90] (-4.85,2.55);
\node[blue!70!black] at (-4.70,3.30) {$\varphi$};

\draw[map, red!70!black]
  (1.25,4.35) to[out=320,in=90] (4.85,2.55);
\node[red!70!black] at (4.70,3.30) {$\psi$};

\draw[tmap] (-3.75,0.72) -- (3.75,0.72);
\node[orange!80!black, above] at (0,0.82) {$\psi\circ\varphi^{-1}$};

\end{tikzpicture}
\end{center}


\begin{theorem}[Boundary points are well-defined]
Let $M$ be a smooth manifold with boundary, and let $p\in M$. Suppose
\[
  (U,\varphi),\qquad (V,\psi)
\]
are two smooth charts containing $p$. Then
\[
  \varphi(p)\in \partial\mathbb H^n
  \iff
  \psi(p)\in \partial\mathbb H^n.
\]
Equivalently,
\[
  \varphi(p)\in \operatorname{Int}\mathbb H^n
  \iff
  \psi(p)\in \operatorname{Int}\mathbb H^n.
\]
\end{theorem}

\begin{proof}
Consider the transition map
\[
  \psi\circ\varphi^{-1}:
  \varphi(U\cap V)\to \psi(U\cap V).
\]
Since $(U,\varphi)$ and $(V,\psi)$ are smoothly compatible charts, this is a smooth transition map between open subsets of $\mathbb H^n$.

Moreover, it restricts to a smooth map
\[
  \psi\circ\varphi^{-1}:
  \partial\mathbb H^n\cap \varphi(U\cap V)
  \to
  \partial\mathbb H^n\cap \psi(U\cap V).
\]
Thus the transition map sends boundary points of $\mathbb H^n$ to boundary points of $\mathbb H^n$.

Therefore,
\[
  \varphi(p)\in \partial\mathbb H^n
  \iff
  \psi(p)\in \partial\mathbb H^n.
\]
Similarly,
\[
  \varphi(p)\in \operatorname{Int}\mathbb H^n
  \iff
  \psi(p)\in \operatorname{Int}\mathbb H^n.
\]
Hence whether $p$ is a boundary point is independent of the chart.
\end{proof}

\begin{definition}
Let $M$ be a smooth manifold with boundary. The \vocab{interior} of $M$ is
\[
  \operatorname{Int} M
  :=
  \{p\in M\mid p \text{ is an interior point}\}.
\]
The \vocab{boundary} of $M$ is
\[
  \partial M
  :=
  \{p\in M\mid p \text{ is a boundary point}\}.
\]
\end{definition}

Thus
\[
  M=\operatorname{Int}M\sqcup \partial M.
\]

\begin{remark}
The interior of $M$ is a smooth $n$-manifold without boundary:
\[
  \operatorname{Int}M
  \quad
  \text{is a smooth $n$-manifold without boundary.}
\]
The boundary of $M$ is a smooth $(n-1)$-manifold without boundary:
\[
  \partial M
  \quad
  \text{is a smooth $(n-1)$-manifold without boundary.}
\]
\end{remark}

\begin{remark}
The topological interior and the manifold interior need not be the same notion.
Similarly, the topological boundary and the manifold boundary need not agree.
\end{remark}

\begin{example}
Let
\[
  M=\{x\in\RR^n\mid x^n>0\}\subset \RR^n.
\]
Then $M$ has no manifold boundary:
\[
  \partial M=\varnothing.
\]
However, as a subset of $\RR^n$, it has topological boundary
\[
  \partial_{\mathrm{top}}M=\RR^{n-1}\times\{0\}.
\]
\end{example}

\begin{example}
If $M$ is an $n$-manifold with boundary, then
\[
  \partial_{\mathrm{top}}M=\varnothing
\]
when $M$ is regarded as a topological space in itself. Indeed,
\[
  \operatorname{Int}_{\mathrm{top}}M=M.
\]
But the manifold boundary $\partial M$ may be nonempty.
\end{example}

\begin{example}
Let
\[
  M=S^n\subset \RR^{n+1}.
\]
Then
\[
  \partial_{\mathrm{mf}}M=\varnothing,
\]
and
\[
  \partial_{\mathrm{top}}M=S^n
\]
if $S^n$ is viewed as a subset of $\RR^{n+1}$.
\end{example}

\begin{example}
Let
\[
  M=\mathbb H^n\cap B(0,1).
\]
Then $M$ is a smooth manifold with boundary. Its manifold boundary is
\[
  \partial M
  =
  \{x\in B(0,1)\mid x^n=0\}.
\]
Geometrically, this is the flat part of the boundary.
\end{example}


\subsection{Smoothness on Half-Spaces}

\begin{remark}
The notion of smoothness for maps defined on subsets of $\mathbb H^n$ is the following.

All partial derivatives should exist and be continuous up to the boundary.
\end{remark}

\begin{definition}
Let
\[
  A\subset \mathbb H^n
\]
be an arbitrary subset, not necessarily open. A map
\[
  f:A\to \RR^m
\]
is called \vocab{smooth} if it can be extended to a smooth map
\[
  \widetilde f:V\to \RR^m
\]
where
\[
  A\subset V\subset \RR^n
\]
and $V$ is open, such that
\[
  \widetilde f|_A=f.
\]
\end{definition}

\begin{center}
\begin{tikzpicture}[
  >=Stealth,
  every node/.style={font=\small},
  outline/.style={thick},
  map/.style={->, thick},
  scale=1
]

% Ambient open set V
\draw[blue!60, thick]
  (0,0)
  .. controls (-1.4,0.9) and (-1.1,2.4) .. (0.4,2.7)
  .. controls (1.8,3.0) and (3.0,2.4) .. (3.3,1.2)
  .. controls (3.6,0.0) and (2.6,-0.8) .. (1.2,-0.7)
  .. controls (0.4,-0.6) and (0.5,-0.1) .. (0,0);

\node[blue!70!black] at (2.6,2.35) {$V$};

% Subset A
\draw[outline]
  (0.6,0.65) -- (1.1,1.15) -- (1.05,0.55)
  -- (1.65,1.15) -- (1.70,0.55)
  -- (2.35,1.20) -- (2.45,0.82);

\node at (1.65,0.30) {$A$};

% Arrow to R^m
\draw[map] (3.8,1.0) -- (6.0,1.0);
\node[above] at (4.9,1.0) {$f$};

\node at (6.6,1.0) {$\RR^m$};

% Extension label
\draw[blue!60, map] (2.1,-0.35) to[out=-20,in=210] (6.0,0.55);
\node[blue!70!black] at (4.2,0.05) {$\widetilde f$};

\end{tikzpicture}
\end{center}

\begin{theorem}[Seeley's Extension Theorem]
Let
\[
  U\subset \mathbb H^n
\]
be open in the relative topology, and let
\[
  f:U\to \RR^m
\]
be a map. Suppose that, for every multi-index $I$, the partial derivative
\[
  \partial^I f
\]
exists on $U\cap \operatorname{Int}\mathbb H^n$ and extends continuously to all of $U$.
Then there exist an open set
\[
  \widetilde U\subset \RR^n,
  \qquad
  U\subset \widetilde U,
\]
and a smooth map
\[
  \widetilde f:\widetilde U\to \RR^m
\]
such that
\[
  \widetilde f|_U=f.
\]

In other words, a map on an open subset of $\mathbb H^n$ is smooth up to the boundary exactly when it can be extended smoothly across the boundary hyperplane $\partial\mathbb H^n$.
\end{theorem}

\begin{definition}[Partial derivatives and multi-index notation]
Let
\[
  W\subset \RR^n
\]
be open, and let
\[
  f:W\to \RR^m.
\]
For each coordinate direction $e_i$, the \vocab{partial derivative} of $f$ in the $x^i$-direction is defined by
\[
  \frac{\partial f}{\partial x^i}(x)
  :=
  \lim_{h\to 0}
  \frac{f(x+h e_i)-f(x)}{h},
\]
provided the limit exists.

A \vocab{multi-index} is an $n$-tuple
\[
  I=(i_1,\dots,i_n),
  \qquad
  i_1,\dots,i_n\in \ZZ_{\ge 0}.
\]
Its order is
\[
  |I|=i_1+\cdots+i_n.
\]
We write
\[
  \partial^I f
  :=
  \frac{\partial^{|I|} f}
  {\partial (x^1)^{i_1}\cdots \partial (x^n)^{i_n}}.
\]
For example, if $n=3$ and $I=(2,0,1)$, then
\[
  \partial^I f
  =
  \frac{\partial^3 f}{\partial (x^1)^2\partial x^3}.
\]
\end{definition}

\begin{remark}
The condition that $\partial^I f$ extends continuously to $U$ means the following. First compute $\partial^I f$ on the interior part
\[
  U\cap \operatorname{Int}\mathbb H^n.
\]
Then, for every boundary point $p\in U\cap \partial\mathbb H^n$, the values of $\partial^I f(x)$ must have a limit as $x$ approaches $p$ from inside $\mathbb H^n$. This limiting value is declared to be the value of the continuous extension at $p$.
\end{remark}

\subsection{Smooth Maps Between Manifolds}

\begin{definition}
Let $M$ be a smooth manifold, possibly with boundary, and let $N$ be a smooth manifold.
A map
\[
  f:M\to N
\]
is called \vocab{smooth} if for every point $p\in M$, there exist smooth charts
\[
  (U,\varphi)
  \quad\text{of }M,
  \qquad
  (V,\psi)
  \quad\text{of }N,
\]
such that
\[
  p\in U,
  \qquad
  f(U)\subset V,
\]
and the coordinate representation
\[
  \widehat f
  :=
  \psi\circ f\circ\varphi^{-1}
  :
  \varphi(U)\to \psi(V)\subset \RR^m
\]
is smooth.
\end{definition}

The map
\[
  \widehat f
  =
  \psi\circ f\circ\varphi^{-1}
\]
is called the \vocab{coordinate representation} of $f$ with respect to the charts
\[
  (U,\varphi)
  \quad\text{and}\quad
  (V,\psi).
\]

\begin{center}
\begin{tikzpicture}[
  >=Stealth,
  every node/.style={font=\small},
  map/.style={->, thick},
  chartmap/.style={->, thick, blue!70!black},
  coordmap/.style={->, thick, blue!70!black},
  outline/.style={thick},
  bluebox/.style={draw=blue!70!black, thick},
  targetbox/.style={draw=black, thick},
  scale=1
]

% Source manifold M
\draw[outline] (-4.8,3.2) ellipse (2.4 and 1.25);
\node at (-6.25,3.85) {$M$};

% U in M
\draw[blue!70!black, thick] (-3.9,3.25) ellipse (0.85 and 0.55);
\node[blue!70!black] at (-3.9,3.42) {$U$};

\fill (-3.9,3.15) circle (1.4pt);
\node[below] at (-3.9,3.08) {$p$};

% Target coordinate image in R^m
\draw[targetbox] (1.8,2.35) rectangle (4.55,4.20);
\node at (3.15,3.95) {$\RR^m$};

\draw[red!70!black, thick] (3.15,3.15) ellipse (0.85 and 0.45);
\node[red!70!black] at (3.15,3.32) {$\psi(V)$};

\fill (3.15,3.05) circle (1.4pt);
\node[below] at (3.15,2.98) {$\psi(f(p))$};

% f arrow
\draw[map] (-2.25,3.25) -- (1.55,3.25);
\node[above] at (-0.35,3.25) {$f$};

% Coordinate image phi(U)
\draw[bluebox] (-2.35,0.05) rectangle (-0.35,1.65);
\node[blue!70!black] at (-0.10,1.35) {$\RR^n$};

\draw[blue!70!black, thick] (-1.35,0.80) ellipse (0.65 and 0.38);
\node[blue!70!black] at (-1.35,0.96) {$\varphi(U)$};

\fill (-1.35,0.72) circle (1.4pt);
\node[below] at (-1.35,0.65) {$\varphi(p)$};

% phi arrow
\draw[chartmap] (-3.55,2.78) -- (-1.70,1.78);
\node[blue!70!black] at (-2.75,2.15) {$\varphi$};

% coordinate representation arrow
\draw[coordmap] (-0.35,1.10) to[out=20,in=225] (2.35,2.35);

\node[blue!70!black] at (1.05,1.85)
  {$\widehat f=\psi\circ f\circ\varphi^{-1}$};

\node[blue!70!black] at (2.25,1.58)
  {coordinate representation};

\end{tikzpicture}
\end{center}

Thus
\[
  C^\infty(M,N)
  :=
  \{f:M\to N\mid f \text{ is smooth}\}.
\]

\begin{lemma}
If
\[
  f:M\to N
\]
is smooth, then for any smooth charts
\[
  (U,\varphi)
  \quad\text{of }M,
  \qquad
  (V,\psi)
  \quad\text{of }N,
\]
the coordinate representation
\[
  \psi\circ f\circ\varphi^{-1}
  :
  \varphi\bigl(U\cap f^{-1}(V)\bigr)
  \to
  \psi(V)
\]
is smooth.
\end{lemma}

\begin{proof}
Let
\[
  p\in U\cap f^{-1}(V).
\]
Since $f$ is smooth, there exist smooth charts
\[
  (\widetilde U,\widetilde\varphi)
  \quad\text{of }M,
  \qquad
  (\widetilde V,\widetilde\psi)
  \quad\text{of }N,
\]
such that
\[
  p\in \widetilde U,
  \qquad
  F(\widetilde U)\subset \widetilde V,
\]
and
\[
  \widetilde\psi\circ F\circ \widetilde\varphi^{-1}
  :
  \widetilde\varphi(\widetilde U)
  \to
  \widetilde\psi(\widetilde V)
\]
is smooth.

On the domain where all expressions are defined, we have
\[
  \psi\circ F\circ \varphi^{-1}
  =
  \left(\psi\circ \widetilde\psi^{-1}\right)
  \circ
  \left(\widetilde\psi\circ F\circ \widetilde\varphi^{-1}\right)
  \circ
  \left(\widetilde\varphi\circ \varphi^{-1}\right).
\]
Here
\[
  \widetilde\varphi\circ \varphi^{-1}
\]
and
\[
  \psi\circ \widetilde\psi^{-1}
\]
are smooth transition maps, while
\[
  \widetilde\psi\circ F\circ \widetilde\varphi^{-1}
\]
is smooth by assumption. Therefore
\[
  \psi\circ F\circ \varphi^{-1}
\]
is smooth near $\varphi(p)$.

Since $p$ was arbitrary, the coordinate representation
\[
  \psi\circ F\circ \varphi^{-1}
  :
  \varphi\bigl(U\cap F^{-1}(V)\bigr)
  \to
  \psi(V)
\]
is smooth.
\end{proof}

\subsection{Basic properties and definitions}

\begin{example}
Let
\[
  f:M\to \RR
\]
be a smooth function.

A smooth function can describe, for example, a ``temperature distribution'' on $M$,
or a ``smooth curve'' if
\[
  f:[0,1]\to M.
\]
\end{example}

\begin{remark}
Read the basic properties and definitions of smooth maps.
\end{remark}

\subsection{Partitions of Unity}

\begin{example}
Given
\[
  f_-,f_+\in C^\infty(\RR),
\]
we want to find
\[
  f\in C^\infty(\RR)
\]
such that
\[
  f=f_-
  \quad\text{on }(-\infty,-1),
\]
and
\[
  f=f_+
  \quad\text{on }(1,\infty).
\]
Choose a cutoff function
\[
  \psi_-\in C^\infty(\RR)
\]
such that
\[
  \psi_-=1
  \quad\text{on }(-\infty,-1],
\]
and
\[
  \operatorname{supp}\psi_-\subset (-\infty,1).
\]
Set
\[
  \psi_+:=1-\psi_-.
\]
Then define
\[
  f:=\psi_-f_-+\psi_+f_+.
\]
Informally,
\[
  \psi_-+\psi_+\equiv 1,
\]
with
\[
  \operatorname{supp}\psi_-\subset (-\infty,1),
  \qquad
  \operatorname{supp}\psi_+\subset (-1,\infty).
\]
This is a simple instance of a partition of unity subordinate to the cover
\[
  \{(-\infty,1),(-1,\infty)\}.
\]
\end{example}

\begin{definition}
Let
\[
  \mathcal X=\{X_\alpha\}_{\alpha\in A}
\]
be an open cover of a topological space $X$. A \vocab{partition of unity subordinate to $\mathcal X$}
is a family of continuous functions
\[
  \{\psi_\alpha:X\to \RR\}_{\alpha\in A}
\]
such that:
\begin{enumerate}
  \item
  \[
    0\le \psi_\alpha\le 1
    \qquad\text{for all }\alpha\in A.
  \]

  \item
  \[
    \operatorname{supp}\psi_\alpha\subset X_\alpha.
  \]

  \item The family
  \[
    \{\operatorname{supp}\psi_\alpha\}_{\alpha\in A}
  \]
  is locally finite.

  \item
  \[
    \sum_{\alpha\in A}\psi_\alpha(x)\equiv 1
    \qquad\text{for all }x\in X.
  \]
\end{enumerate}
\end{definition}

\begin{remark}
The local finiteness condition implies that
\[
  \sum_{\alpha\in A}\psi_\alpha
\]
is locally a finite sum. Since
\[
  \sum_{\alpha\in A}\psi_\alpha\equiv 1,
\]
we also have
\[
  X=\bigcup_{\alpha\in A}\operatorname{supp}\psi_\alpha.
\]
\end{remark}

\begin{theorem}[Existence of smooth partitions of unity]
Let $M$ be a smooth manifold, possibly with boundary. For any open cover
\[
  \mathcal X=\{X_\alpha\}_{\alpha\in A}
\]
of $M$, there exists a smooth partition of unity
\[
  \{\psi_\alpha\}_{\alpha\in A}
\]
subordinate to $\mathcal X$.
\end{theorem}

\begin{remark}
Equivalently, for any open cover of a smooth manifold, one can find smooth bump functions whose supports lie inside the elements of the cover and whose locally finite sum is identically equal to $1$.
\end{remark}
